\documentclass[11pt,a4paper]{article}

%% PACHETE MATEMATICE
\usepackage{amsmath,amssymb,amsfonts}
\usepackage{amsthm}
\usepackage{mathpazo}
\usepackage{eulervm}
\usepackage{enumerate}
\usepackage{color}
\usepackage{graphicx}
\usepackage[all]{xy}
\usepackage{xcolor}
\usepackage{framed}
\usepackage[unicode]{hyperref}
\setcounter{MaxMatrixCols}{20}
\usepackage{multicol}
%\usepackage[romanian]{babel}


%% ASPECT
\linespread{1.05}
\usepackage[T1]{fontenc}
\usepackage[utf8x]{inputenc}
\usepackage[margin=2cm]{geometry}
% \usepackage[color]{showkeys}
\usepackage{anyfontsize}
\usepackage{titlesec}
\titleformat*{\section}{\large\bfseries}

\usepackage{fancyhdr}
\pagestyle{fancy}
\rhead{\small \color{gray} Notițe de Adrian Manea}
\lhead{\small \color{gray} BCD-Aero, 2017---2018, L. Costache \& M. Olteanu}


\usepackage[autostyle = false, style = german]{csquotes}
\usepackage[romanian]{babel}
\newcommand\qq{\enquote}

%% COMENZILE MELE
\renewcommand*{\proofname}{Dem.:}
\newcommand\bb{\mathbb}
\newcommand\dr{\mathrm}
\newcommand\kal{\mathcal}
\newcommand\fk{\mathfrak}
\newcommand\op{\oplus}
\newcommand\ot{\otimes}
\newcommand\ds{\displaystyle}
\newcommand\id{\indent}
\newcommand\nid{\noindent}
\newcommand\seq{\subseteq}
\newcommand\sneq{\subsetneq}
\newcommand\speq{\supseteq}
\newcommand\spneq{\supsetneq}
\newcommand\rar{\rightarrow}
\newcommand\Rar{\Rightarrow}
\newcommand\xrar{\xrightarrow}
\newcommand\lar{\leftarrow}
\newcommand\Lar{\Leftarrow}
\newcommand\xlar{\xleftarrow}
\newcommand\lrar{\leftrightarrow}
\newcommand\Lrar{\Leftrightarrow}
\newcommand\ideal{\trianglelefteq}
\newcommand\ai{\text{ a.\^{i}. }}
\newcommand\sk{(\textit{Schi\c{t}\u{a}}) }
\newcommand\rhup{\rightharpoonup}
\newcommand\rhdn{\rightharpoondown}
\newcommand\lhup{\leftharpoonup}
\newcommand\lhdn{\leftharpoondown}
\newcommand\vid{\emptyset}
\newcommand\wh{\widehat}
\newcommand\ol{\overline}
\newcommand\NN{\mathbb{N}}
\newcommand\RR{\mathbb{R}}
\newcommand\ZZ{\mathbb{Z}}
\newcommand\QQ{\mathbb{Q}}
\newcommand\CC{\mathbb{C}}

%% STILURILE MELE
\newtheoremstyle{dotless}{}{}{\itshape}{}{\bfseries}{:}{ }{}

\theoremstyle{dotless}
\newtheorem{theorem}{Teorem\u{a}}[section]
\newtheorem{proposition}{Propozi\c{t}ie}[section]
\newtheorem{lemma}{Lem\u{a}}[section]
\newtheorem{corollary}{Corolar}[section]
\newtheorem{exercise}{Exerci\c{t}iu}

\newtheoremstyle{dotlessdr}{}{}{}{}{\bfseries}{:}{ }{}
\theoremstyle{dotlessdr}
\newtheorem{example}{Exemplu}[section]
\newtheorem{definition}{Defini\c{t}ie}[section]
\newtheorem{remark}{Observa\c{t}ie}[section]




\begin{document}

\begin{center}
  {\large\textbf{Seminar 12 \\
    Probleme de extrem}}
\end{center}

1. Să se calculeze extremele locale ale funcțiilor:
\begin{enumerate}[(a)]
\item $f : \RR^2 \to \RR, f(x, y) = x^3 + y^3 - 6xy$;
\item $g : \RR^2 \to \RR, g(x, y) = x^3 + 8y^3 - 2xy$;
\item $h : \RR^2 \to \RR, h(x, y) = x^2ye^{2x + 3y}$;
\item $k : \RR^2 \to \RR, k(x, y) = xy$;
\end{enumerate}

\emph{Soluție (schiță)} (c): Mulțimea punctelor critice este:
\[
  \{(0, y) \mid y \in \RR \} \cup \{(-1, -\frac{1}{3}) \}.
\]

Punctul izolat este punct de minim. În punctele $(0, y)$, avem o valoare proprie
nulă pentru hessiană. Se evaluează, atunci, semnul diferenței:
\[
  f(x, y) - f(0, y) = x^2ye^{2x+3y},
\]
care ar trebui să aibă semn constant. Remarcăm, însă, că pentru punctele situate
deasupra axei $OX$, adică cu $y > 0$, există un disc de rază mică (de exemplu, $\frac{y}{2}$),
pe are diferența $f(x, y) - f(0, y)$ este pozitivă, deci aceste puncte sînt minime
locale. Analog, punctele de sub axa $OX$ sînt maxime locale. Pentru $y = 0$, adică
în origine, diferența $f(x, y)$ nu păstrează semn constant pe nicio vecinătate a
lui $(0, 0)$, deci nu este punct de extrem.

\vspace{1cm}

2. Să se determine punctele de extrem ale funcției implicite $y = f(x)$, definite prin
ecuația $F(x,y) = x^2 - 2xy + 5y^2 - 2x + 4y + 1 = 0$.

\vspace{1cm}

3. Să se determine punctele de extrem și valorile extreme pentru funcțiile:
\begin{enumerate}[(a)]
\item $f : \RR^2 \to \RR, f(x, y) = (x + y^2) \cdot e^{x + 2y}$;
\item $f : \RR^2 \to \RR, f(x, y) = e^{x^2 - y} + e^{x + y}$;
\item $f : \RR^2 \to \RR, f(x, y) = \sin x \sin y \sin (x + y)$;
\item $f : (\RR - \{0\})^2 \to \RR, f(x, y) = \frac{x^2}{2y} + y^2 + \frac{8}{x} + 5$.
\end{enumerate}

\vspace{1cm}

4. Determinați valorile extreme pentru funcțiile $f$ definite pe mulțimile $K$,
în cazurile:
\begin{enumerate}[(a)]
\item $f(x, y) = 5x^2 + 4xy + 8y^2, K = \{(x,y) \in \RR^2 \mid x^2 + y^2 \leq 1 \}$;
\item $f(x, y) = x + y - 2xy, K_1 = \{ (x, y) \in \RR^2 \mid x^2 + y^2 \leq 1 \}$,
  precum și $K_2 = [0, 1] \times [0, 1]$;
\item $f(x, y) = xy(2x + 3y - 5), K = \{(x, y) \in RR^2 \mid 2x + 3y \leq 12, x,y \geq 0\}$;
\item $f(x, y, z) = x^2 + y^2 + xz - yz, K: x + y + z = 1, x,y,z \geq 0$;
\item $f(x, y, z) = \cos x + \cos y + \cos z, K: x + y + z = \pi, x, y, z \geq 0$;
\item $f(x, y, z) = 2x^2 + 2y^2 + z^2 + 2xy + 2xz - 8x - 6y - 4z$, unde
  $K: x + y + z = 3, x,y,z \geq 0$;
\item $f(x,y) = x^2 + y^2 - x - y, K: x + y = 1$.
\end{enumerate}

\vspace{1cm}

5. Să se determine punctele cele mai depărtate de origine care se găsesc pe suprafața
de ecuație $4x^2 + y^2 + z^2 - 8x - 4y + + 4 = 0$.

\vspace{1cm}

6. Să se determine valorile minime ale produsului $xy$, cînd $x$ și $y$ sînt coordonatele
unui punct de pe elipsa de ecuație $x^2 + 2y^2 = 1$.


\end{document}